\documentclass[../../main.tex]{subfiles}
\graphicspath{{\subfix{../../image/}}}

\begin{document}

\section*{第三章小结}\label{section:ch03-第三章小结}

本章中我们详细地研究了随机变量。用随机变量描述随机现象是近代概率论中最重要的方法，以后我们所讨论的随机事件几乎都用随机变量来描述。

我们给出了随机变量的严格定义，按照这种观点，随机变量是定义在样本空间上的具有某种可测性的实值函数，只是出于历史的原因，才沿用“变量”二字。

对于随机变量，重要的是要知道它取哪一些值以及以怎样的概率取这些值。从这个角度讲，分布函数完整地描述了随机变量。同时，分布函数具有良好的分析性质，便于研究，因此它成为研究随机变量的重要工具。

离散型随机变量与连续型随机变量是最重要的两类随机变量，由于它们的取值特点不同，因此对它们的描述及处理方法都有很大不同。前者列出取值及相应的概率分布，所用的数学工具主要是求和与级数；后者则用密度函数，广泛使用微积分。应该进行对比，从而加深理解。二者的统一是分布函数，它描述了一切随机变量。

一种分布提供一个数学模型。分布函数是概率论的理论与应用的重要结合点。概率论中有名的分布函数大多在本章的正文、例子或习题中出现，数理统计中的三大分布——$\chi^{2}$ 分布、$F$ 分布、$t$ 分布也被导出。我们相当注意揭示各种分布函数的特征性质及它们之间的联系。这些分布函数在不同的理论和实际问题中扮演重要角色。

正态分布是概率论中最重要的分布，在应用中及理论研究中占有头等重要的地位，它与 Poisson 分布及二项分布是概率论中最重要的三种分布。判断一种分布重要性的标准是：(1) 在实际工作中经常遇到；(2) 在理论研究中重要，具有较好的性质；(3) 用它能导出许多重要分布。以上三种分布都满足这些要求。

把几个随机变量放在一起作为随机向量研究时不但需要研究各个分量个别的性质，而且要考虑它们之间的联系，从而大大丰富了研究的内容，条件分布及独立性概念也随之出现，它们是条件概率及事件独立性概念在随机变量场合的具体化，在今后研究中很重要。

随机变量的函数的分布律的推导，在数理统计中及在概率论的许多应用中相当重要，我们这里分一对一，多对一，多对多三种类型对直接法和变换法这两种处理方法作了深入介绍。这部分内容，只有通过动手多作练习，才能牢固掌握。

可测性是严格定义随机变量的关键，因而离不开 Borel 点集与 Borel 函数等概念，否则连随机变量的函数是否随机变量都无法讲清。

\end{document}